% ============================================================================
% "Geometry and topology of discrete Maxwell wave speed"
% Companion to JCP Paper 1 (JCOMP-D-26-00537)
% Target: J. Comput. Phys.
% ============================================================================
\documentclass[preprint,12pt]{elsarticle}
\usepackage{amsmath,amssymb,amsthm}
\usepackage{bm}
\usepackage{mathtools}
\usepackage{booktabs}
\usepackage{graphicx}

\newtheorem{theorem}{Theorem}[section]
\newtheorem{lemma}[theorem]{Lemma}
\newtheorem{proposition}[theorem]{Proposition}
\newtheorem{corollary}[theorem]{Corollary}
\newtheorem{remark}{Remark}[section]

\DeclareMathOperator{\Vol}{Vol}
\DeclareMathOperator{\tr}{tr}
\DeclareMathOperator{\rank}{rank}
\DeclareMathOperator{\im}{im}
\DeclareMathOperator{\Span}{span}

\begin{document}

\begin{frontmatter}
\title{Geometry and topology of discrete Maxwell wave speed}
\author{Alexandru Toader}

\begin{abstract}
We prove that on any periodic Voronoi tessellation in three dimensions,
equipped with an exact discrete cochain complex ($d_1 d_0 = 0$), the
long-wavelength electromagnetic wave speed is exactly unity in natural
units. The result $\omega^2 = |k|^2 + O(|k|^4)$ is a geometric identity
at finite resolution, not a convergence statement. The proof rests on two
conditions: (A)~the Voronoi metric tensors satisfy
$G = H = \Vol \cdot I$ (proved via the divergence theorem), and (B)~the
Bloch-periodic cochain complex preserves exactness (provided by a companion
paper). These combine through a Schur complement reduction and a discrete
curl identity to yield the exact leading dispersion. Numerical verification
spans three ordered foams (Kelvin, Weaire--Phelan, C15) and five random
Voronoi tessellations. Removing either condition produces $O(1)$ error in
the wave speed, demonstrating that both geometry and topology are essential.
\end{abstract}

\begin{keyword}
discrete exterior calculus \sep Maxwell equations \sep Voronoi tessellation
\sep wave speed \sep cochain complex
\end{keyword}
\end{frontmatter}


% ============================================================================
% §1. INTRODUCTION
% ============================================================================
\section{Introduction}
\label{sec:intro}

Discrete exterior calculus (DEC) provides a coordinate-free framework for
discretizing differential forms on polyhedral
complexes~\cite{hirani2003,desbrun2005,arnold2006}. For Maxwell's equations
on a periodic complex, the discrete curl--curl operator
$K = d_1^\dagger \star_2 \, d_1$ with mass matrix $M = \star_1$ defines a
generalized eigenvalue problem
\begin{equation}
  K \psi = \omega^2 M \psi
  \label{eq:eigenvalue}
\end{equation}
at Bloch wavevector~$k$. The wave speed~$c$, defined by
$\lim_{|k|\to 0} \omega^2/|k|^2$, is a fundamental physical quantity that
should equal unity in natural units.

On structured grids, $c = 1$ follows from lattice symmetry. On unstructured
polyhedral meshes, the wave speed depends on the Hodge star
operators~$\star_1$, $\star_2$ (encoding metric information) and on the
coboundary operator~$d_1$ (encoding topology). Whether an unstructured mesh
can achieve $c = 1$ without parameter fitting has remained an open question.

In this paper we prove that $c^2 = 1$ on any periodic Voronoi tessellation
equipped with an exact cochain complex. Two conditions are sufficient:
\begin{itemize}
\item[\textbf{(A)}] \textbf{Metric.} The edge and face Gram tensors satisfy
  $G = H = \Vol \cdot I$, where
  $G_{ij} = \sum_e \star_1[e]\, (\Delta x_e)_i (\Delta x_e)_j$ and
  $H_{ij} = \sum_f \star_2[f]\, (A_f)_i (A_f)_j$.
\item[\textbf{(B)}] \textbf{Topology.} The Bloch-periodic cochain complex
  preserves exactness: $d_1(k)\, d_0(k) = 0$ for all wavevectors~$k$.
\end{itemize}
Condition~(A) is a geometric identity proved via the divergence theorem on
Voronoi and dual cells (\S\ref{sec:metric}). Condition~(B) is provided by
the exactness-preserving construction of~\cite{paper1}. Necessity of~(A) is
proved analytically: $c = 1$ for all directions and polarizations implies
$G = H = \lambda I$ (\S\ref{sec:converse}). Necessity of~(B) is supported
numerically: all tested violations of exactness yield $c \neq 1$ with $O(1)$
error (\S\ref{sec:exactness}).

The theorem states $\omega^2 = |k|^2 + O(|k|^4)$: the leading dispersion
coefficient is exactly unity at any finite mesh resolution. This is a
geometric identity, not a convergence statement---no mesh refinement or
effective medium limit is required. The $O(|k|^4)$ dispersion correction is
structure-dependent (analogous to $\omega = 2\sin(k/2)$ on a 1D chain), but
the leading term is universal across all periodic Voronoi tessellations. The
Voronoi DEC with exact complex is ``already homogenized'': the discrete
metric tensors are proportional to the identity, the discrete analog of a
trivially isotropic medium.

This paper is companion to~\cite{paper1}, which constructs the exact cochain
complex (condition~B) on periodic polyhedral meshes. Here we prove the
physical consequence: the discrete wave speed equals the continuum value.
Together, the two papers provide a complete discrete Maxwell theory on
periodic Voronoi complexes with exact topology and exact long-wavelength
wave speed.

The paper is organized as follows. Section~\ref{sec:metric} proves the
metric identities $G = H = \Vol \cdot I$.
Section~\ref{sec:theorem} states and proves the main theorem $c^2 = 1$ via
Schur complement reduction and a discrete curl identity, followed by the
converse in \S\ref{sec:converse}.
Sections~\ref{sec:exactness} and~\ref{sec:geometry} demonstrate the
necessity of both conditions by removing each independently.
Section~\ref{sec:spectral} analyzes the spectral structure of the
approximation error. Section~\ref{sec:discussion} discusses implications
and extensions. Appendix~\ref{app:naive_pt} explains why naive perturbation
theory fails, and Appendix~\ref{app:std_anatomy} provides a detailed anatomy
of the standard DEC failure mode.


% --- 1.1 Exact periodic cochain complex ---
\subsection{Exact periodic cochain complex (condition~B)}
\label{sec:exact_complex}

The standard Bloch-periodic coboundary operators are obtained by
multiplying each incidence entry by a phase $e^{ik \cdot n_e L}$, where
$n_e \in \mathbb{Z}^3$ is the lattice shift of edge~$e$. This simple
phasing does not preserve the cochain complex: generically
$d_1(k)\, d_0(k) \neq 0$ at $k \neq 0$.

The construction of~\cite{paper1} modifies the $1$-coboundary operator
via a recursive correction that redistributes lattice shifts among the
edges of each face. The corrected operator $d_{1,\mathrm{ex}}(k)$
satisfies $d_1(k)\, d_0(k) = 0$ for all~$k$ while preserving the
$\Gamma$-point operator $d_1(0)$ and the support structure of the
incidence matrix. The $2$-coboundary $d_2(k)$ requires no correction:
$d_2 d_1 = 0$ holds automatically (it is the topological identity
$\partial^2 = 0$). Full details are given in~\cite{paper1}.


% ============================================================================
% §2. METRIC IDENTITIES
% ============================================================================
\section{Metric identities: $G = H = \Vol \cdot I$}
\label{sec:metric}

Consider a periodic Voronoi tessellation of the 3-torus~$\mathbb{T}^3$ with
total volume~$\Vol$. The CW complex has vertex set~$V$ (Voronoi vertices,
where four or more cells meet), edge set~$E$ (Voronoi edges), face set~$F$
(Voronoi faces), and cell set~$C$ (Voronoi cells, one per seed). Each edge
$e = (v_1, v_2)$ carries a vector $\Delta x_e = x_{v_2} - x_{v_1}$
(modulo periodicity) and each face~$f$ carries an oriented area
vector~$A_f$.

The diagonal Hodge star operators are defined by primal--dual
perpendicularity:
\begin{equation}
  \star_1[e] = \frac{|f^*_e|}{|\Delta x_e|}, \qquad
  \star_2[f] = \frac{|e^*_f|}{|A_f|},
  \label{eq:hodge_stars}
\end{equation}
where $f^*_e$ is the dual face of edge~$e$ (the polygon of cell centers
around~$e$, with area~$|f^*_e|$) and $e^*_f$ is the dual edge of face~$f$
(the segment connecting the two cell centers separated by~$f$, with
length~$|e^*_f|$).

We define the \emph{edge metric tensor}~$G$ and \emph{face metric tensor}~$H$:
\begin{equation}
  G_{ij} = \sum_{e \in E} \star_1[e]\, (\Delta x_e)_i\, (\Delta x_e)_j,
  \qquad
  H_{ij} = \sum_{f \in F} \star_2[f]\, (A_f)_i\, (A_f)_j.
  \label{eq:GH_def}
\end{equation}

% --- 2.1 Edge tensor G ---
\subsection{Edge metric: $G = \Vol \cdot I$}
\label{sec:G_proof}

\begin{theorem}[Edge metric identity]
\label{thm:G}
On any periodic Voronoi tessellation of\/ $\mathbb{T}^3$,
$G = \Vol \cdot I$.
\end{theorem}

\begin{proof}
The dual of the Voronoi complex is the Delaunay complex. Each Voronoi
vertex~$v$ is surrounded by a dual 3-cell~$D_v$ (the convex hull of the
cell centers adjacent to~$v$). We apply the divergence theorem to the
vector field $F = x_i\, \hat{e}_j$ on each dual cell:
\begin{equation}
  \Vol(D_v)\, \delta_{ij}
    = \sum_{\substack{e \;\text{at}\; v}} \int_{f^*_e}
      x_i\, (n_e)_j \, dA,
  \label{eq:div_thm_G}
\end{equation}
where the sum runs over all edges~$e$ emanating from~$v$, and $n_e$ is the
outward unit normal to the dual face~$f^*_e$.

\emph{Perpendicularity.}
The dual face~$f^*_e$ is the polygon of cell centers around edge~$e$. Each
cell center lies on a Voronoi face containing~$e$, and every such face is
the perpendicular bisector of a segment between two cell centers. Since $e$
lies in each of these bisector planes, $e$ is perpendicular to every
segment connecting adjacent cell centers around it. The polygon of cell
centers therefore lies in a plane orthogonal to~$\Delta x_e$. Its outward
normal is $n_e = \pm \Delta x_e / |\Delta x_e|$.

\emph{Decomposition.}
On dual face~$f^*_e$, decompose each point as $x = p_e + \delta$, where
$p_e = x_v + t_e\, \Delta x_e$ is the point where the dual face plane
intersects edge~$e$ (at parameter~$t_e$ from vertex~$v$), and
$\delta \perp \Delta x_e$ satisfies $\int_{f^*_e} \delta\, dA
= |f^*_e|\, \eta_e$ with $\eta_e$ the in-plane centroid offset.
Substituting into~\eqref{eq:div_thm_G} yields three terms:
\begin{align}
  \text{Term 1:} &\quad
    (x_v)_i \sum_{e\, \text{at}\, v} \star_1[e]\, (\Delta x_e)_j,
  \label{eq:term1_G} \\
  \text{Term 2:} &\quad
    \sum_{e\, \text{at}\, v} t_e\, \star_1[e]\,
    (\Delta x_e)_i\, (\Delta x_e)_j,
  \label{eq:term2_G} \\
  \text{Term 3:} &\quad
    \sum_{e\, \text{at}\, v} (\eta_e)_i\, \star_1[e]\, (\Delta x_e)_j,
  \label{eq:term3_G}
\end{align}
using $|f^*_e| = \star_1[e] \cdot |\Delta x_e|$ and
$n_e = \Delta x_e / |\Delta x_e|$.

\emph{Term~1 vanishes.}
For a closed dual cell, $\sum_{e\, \text{at}\, v} n_e\, |f^*_e| = 0$
(flux of a constant through a closed surface), which gives
$\sum_{e\, \text{at}\, v} \star_1[e]\, (\Delta x_e)_j = 0$.

\emph{Term~2 yields~$G$.}
Sum~\eqref{eq:div_thm_G} over all vertices~$v$. Each edge $e = (v,w)$
appears twice: from vertex~$v$ with parameter~$t_e^v$ and outward
normal~$+\Delta x_e/|\Delta x_e|$, and from vertex~$w$ with
parameter~$t_e^w = 1 - t_e^v$ and outward normal~$-\Delta x_e/|\Delta x_e|$.
Since $(\Delta x_e)_i(\Delta x_e)_j$ is invariant under sign reversal, the
total contribution of edge~$e$ is
$(t_e^v + t_e^w)\, \star_1[e]\, (\Delta x_e)_i (\Delta x_e)_j
= \star_1[e]\, (\Delta x_e)_i (\Delta x_e)_j$.
Thus $\sum_v \text{Term~2} = G_{ij}$.

\emph{Term~3 cancels pairwise.}
Each dual face~$f^*_e$ is shared by the dual cells at the two endpoints
of~$e$. The centroid offset~$\eta_e$ is intrinsic to the dual face (it
depends only on the positions of the cell centers forming the polygon) and
is the same from both sides. The normals are opposite.
Periodicity ensures every dual face in the fundamental domain has a partner
(boundary faces pair with their periodic images). The contributions cancel:
$\sum_v \text{Term~3} = 0$.

Collecting terms: $\Vol\, \delta_{ij} = G_{ij}$.
\end{proof}

% --- 2.2 Face tensor H ---
\subsection{Face metric: $H = \Vol \cdot I$}
\label{sec:H_proof}

\begin{theorem}[Face metric identity]
\label{thm:H}
On any periodic Voronoi tessellation of\/ $\mathbb{T}^3$,
$H = \Vol \cdot I$.
\end{theorem}

\begin{proof}
Apply the divergence theorem to $F = x_i\, \hat{e}_j$ on each Voronoi
cell~$C$:
\begin{equation}
  \Vol(C)\, \delta_{ij}
    = \sum_{f \in \partial C} \int_f x_i\, (n_f)_j\, dA.
  \label{eq:div_thm_H}
\end{equation}

\emph{Perpendicularity.}
Each Voronoi face~$f$ is the perpendicular bisector of the dual edge~$e^*_f$
connecting the two cell centers it separates. The outward face normal
satisfies $n_f \parallel e^*_f$, i.e.\ $n_f = \pm A_f / |A_f|$.

\emph{Decomposition.}
Since face~$f$ is the perpendicular bisector of~$e^*_f$, the midpoint
$m_f = (p_C + p_{C'})/2$ of the dual edge lies on the face, at distance
$|\Delta p_f|/2$ from~$p_C$ along~$n_f$, where $\Delta p_f = p_{C'} - p_C$.
Decompose each point on face~$f$ as $x = p_C + (m_f - p_C) + \delta$,
where $\delta \perp n_f$ satisfies
$\int_f \delta\, dA = |A_f|\, \eta_f$ with $\eta_f$ the in-plane
centroid offset. As in the $G$~proof, three terms arise:
\begin{align}
  \text{Term 1:} &\quad
    (p_C)_i \sum_{f \in \partial C} (n_f)_j\, |A_f| = 0
    \quad\text{(closed cell)},
  \label{eq:term1_H} \\
  \text{Term 2:} &\quad
    \tfrac{1}{2} \sum_{f \in \partial C}
    \star_2[f]\, (A_f)_i\, (A_f)_j,
  \label{eq:term2_H} \\
  \text{Term 3:} &\quad
    \sum_{f \in \partial C} (\eta_f)_i\, (n_f)_j\, |A_f|.
  \label{eq:term3_H}
\end{align}
In Term~2 we used $(m_f - p_C)_i = (\Delta p_f)_i / 2$ and
$(n_f)_j\, |A_f| = \pm (A_f)_j$ with
$(\Delta p_f)_i (A_f)_j / |\Delta p_f|
= \star_2[f] \cdot (A_f)_i (A_f)_j / |A_f|$
since $\Delta p_f \parallel A_f$.

Summing over all cells: each face is shared by two cells contributing
$\tfrac{1}{2}$ each, so $\sum_C \text{Term~2} = H_{ij}$.
The centroid offset~$\eta_f$ is face-intrinsic (identical from both cells);
opposite normals cancel Term~3 pairwise.
Therefore $\Vol\, \delta_{ij} = H_{ij}$.
\end{proof}

\begin{remark}
The two proofs are structurally identical: apply the divergence theorem
on dual 3-cells for~$G$ and on primal 3-cells for~$H$. In each case,
Voronoi perpendicularity aligns the outward normal with the relevant
direction vector ($\Delta x_e$ or $A_f$), the closed-surface condition
kills Term~1, and pairwise cancellation eliminates Term~3.
\end{remark}

% --- 2.3 Admissible perturbations ---
\subsection{Admissible perturbations}
\label{sec:admissible}

The constraint $G = \Vol \cdot I$ imposes 6 scalar conditions (the
independent components of a $3 \times 3$ symmetric tensor) on the $|E|$
Hodge star values $\star_1[e]$. On all tested structures, the constraint
matrix $(\Delta x_e)_i (\Delta x_e)_j$ has rank~6, so the admissible set
\begin{equation}
  \mathcal{A} = \{ \star_1 \in \mathbb{R}^{|E|}_{>0} :
    G(\star_1) = \Vol \cdot I \}
\end{equation}
is a smooth manifold of dimension $|E| - 6$, intersected with the positive
orthant.

\begin{proposition}
\label{prop:rank6}
If every vertex has degree $\geq 4$ and the edge directions at each vertex
span~$\mathbb{R}^3$, then
$\rank\{(\Delta x_e) \otimes (\Delta x_e) : e \in E\} = 6$,
so $\dim \mathcal{A} = |E| - 6$.
\end{proposition}

\begin{proof}
The map $v \mapsto v \otimes v$ sends $\mathbb{R}^3$ into the 6-dimensional
space of symmetric matrices~$\mathrm{Sym}^2(\mathbb{R}^3)$. Since three
linearly independent vectors $v_1, v_2, v_3$ produce six linearly
independent symmetric matrices $\{v_i \otimes v_i,\,
(v_i + v_j) \otimes (v_i + v_j)\}_{i<j}$ (the latter isolating
cross terms $v_i \otimes v_j + v_j \otimes v_i$), the spanning condition
ensures rank~6.
\end{proof}

The Voronoi Hodge star is one point in~$\mathcal{A}$; the remaining
$|E|-6$ degrees of freedom parameterize perturbations of~$\star_1$ that
preserve $c = 1$. Numerically, the admissible cone
radius ranges from $1.4$ to $6 \times \overline{\star_1}$ across tested
structures, indicating a substantial basin around the Voronoi point.

% --- 2.4 Numerical verification ---
\subsection{Numerical verification}
\label{sec:metric_numerics}

Table~\ref{tab:metric} summarizes the metric identity on all tested
structures. The identity $G = H = \Vol \cdot I$ holds to machine precision
on ordered foams and to $O(10^{-10})$ on random Voronoi tessellations (the
residual reflects Qhull vertex-placement tolerance).

\begin{table}[ht]
\centering
\caption{Metric identity verification. $\epsilon_G = \|G/\Vol - I\|_\infty$,
$\epsilon_H = \|H/\Vol - I\|_\infty$.}
\label{tab:metric}
\begin{tabular}{lrrcc}
\toprule
Structure & $|V|$ & $|E|$ & $\epsilon_G$ & $\epsilon_H$ \\
\midrule
Kelvin ($N\!=\!2$) & 96 & 192 & $10^{-16}$ & $10^{-16}$ \\
C15 ($N\!=\!2$) & 432 & 864 & $10^{-15}$ & $10^{-16}$ \\
Weaire--Phelan ($N\!=\!2$) & 368 & 736 & $10^{-16}$ & $10^{-16}$ \\
Random ($n\!=\!80$, 5~seeds) & 540--570 & 1080--1140
  & $10^{-12}$ & $10^{-12}$ \\
Degenerate ($n\!=\!60$, ratio 40$\times$) & 410 & 820 & $10^{-12}$ & $10^{-12}$ \\
\bottomrule
\end{tabular}
\end{table}

Additional verifications:
the trace identity $\tr G = \tr H = 3\Vol$ holds on all structures;
perpendicularity $\hat{e}^* \parallel \hat{n}_f$ is verified to machine precision
(dual edge direction parallel to face normal);
and domain scaling $G \to s^3 I$ under $V \to sV$, $L \to sL$ confirms
dimensional correctness.

% --- 2.5 Tight frame remark ---
\subsection{Tight frame interpretation}
\label{sec:tight_frame}

\begin{remark}
The weighted edge vectors $f_e = \sqrt{\star_1[e]}\, \Delta x_e$ form a
tight frame for~$\mathbb{R}^3$ with frame bound $A = \Vol$:
\begin{equation}
  F^\dagger F = \Vol \cdot I, \qquad
  \sigma_1 = \sigma_2 = \sigma_3 = \sqrt{\Vol},
\end{equation}
where $F$ is the $|E| \times 3$ matrix with rows $f_e^\top$.
The identity $G = \Vol \cdot I$ is equivalent to Parseval-type isotropic
quadrature: the Hodge star--weighted edge directions sample all of
$\mathbb{R}^3$ uniformly.
\end{remark}


% ============================================================================
% §3. THEOREM: c² = 1
% ============================================================================
\section{Acoustic wave speed}
\label{sec:theorem}

% --- 3.1 Setup and statement ---
\subsection{Setup and statement}
\label{sec:setup}

The discrete Maxwell operator at Bloch wavevector~$k$ is
$K(k) = d_1(k)^\dagger \star_2\, d_1(k)$ acting on 1-forms, with mass
matrix $M = \star_1$. The generalized eigenvalue problem~\eqref{eq:eigenvalue}
determines the discrete dispersion relation $\omega^2(k)$.

At $k = 0$ (the $\Gamma$~point), the first cohomology
$H^1(\mathbb{T}^3) \cong \mathbb{R}^3$ is spanned by harmonic 1-forms
\begin{equation}
  (h_\alpha)_e = \frac{(\Delta x_e)_\alpha}{\sqrt{\Vol}}, \quad \alpha = 1,2,3,
  \label{eq:harm}
\end{equation}
which are $M$-orthonormal:
$h_\alpha^\dagger M\, h_\beta
  = G_{\alpha\beta}/\Vol = \delta_{\alpha\beta}$
(using $G = \Vol \cdot I$ from Theorem~\ref{thm:G}).
At $k \neq 0$, the cohomology $H^1(k)$ vanishes---verified numerically on
all tested structures. (Dimension count: with exactness,
$\dim\ker(d_1) = |V|$ and $\dim\im(d_0) = |V|$ at generic~$k$, so
$H^1 = \ker d_1 / \im d_0 = 0$.)
The three zero modes at~$\Gamma$ deform into one gradient mode and two
acoustic modes.

The second cohomology $H^2(\mathbb{T}^3) \cong \mathbb{R}^3$ at~$\Gamma$
is spanned by
\begin{equation}
  (h_2)_\beta = \frac{(A_f)_\beta}{\sqrt{\Vol}}, \quad \beta = 1,2,3,
  \label{eq:h2}
\end{equation}
which are $\star_2$-orthonormal:
$(h_2)_\beta^\dagger \star_2\, (h_2)_{\beta'}
  = H_{\beta\beta'}/\Vol = \delta_{\beta\beta'}$
(using $H = \Vol \cdot I$ from Theorem~\ref{thm:H}).

\begin{theorem}[Wave speed identity]
\label{thm:c_squared}
On any periodic Voronoi tessellation of\/ $\mathbb{T}^3$ with exact
cochain complex ($d_1(k)\, d_0(k) = 0$ for all~$k$),
\begin{equation}
  \omega^2 = |k|^2 + O(|k|^4),
  \label{eq:c_squared}
\end{equation}
giving $c^2 = \lim_{|k|\to 0} \omega^2/|k|^2 = 1$.
\end{theorem}

% --- 3.2 Proof ---
\subsection{Proof of Theorem~\ref{thm:c_squared}}
\label{sec:proof}

The proof uses four ingredients:
\begin{enumerate}
\item $H = \Vol \cdot I$, so $(h_2)^\dagger \star_2\, (h_2) = I_3$
  \hfill (Theorem~\ref{thm:H})
\item $d_1(0)\, h_\alpha = 0$: the curl of a constant form vanishes.
\item Cochain exactness at all~$k$:
  \begin{itemize}
  \item $d_1(k)\, d_0(k) = 0$ (from~\cite{paper1}; the non-trivial condition),
  \item $d_2(k)\, d_1(k) = 0$ (a topological identity, $\partial^2 = 0$ on
    any oriented cell complex; automatic).
  \end{itemize}
\item Discrete curl identity (Lemma~\ref{lem:curl} below).
\end{enumerate}

\emph{Step 1: image of harmonics.}
Define the $|F| \times 3$ matrix $u = d_1(k)\, h$, where $h$ is the
$|E| \times 3$ matrix of harmonic 1-forms~\eqref{eq:harm}. Since
$d_1(0)\, h = 0$ (ingredient~2), we have
$u = [d_1(k) - d_1(0)]\, h = O(k)$.

\emph{Step 2: orthogonal decomposition in face space.}
Let $W = \im(d_1(0)) \subset \mathbb{C}^{|F|}$ and decompose
\begin{equation}
  u = u_{\mathrm{im}} + u_\perp,
  \label{eq:u_decomp}
\end{equation}
$\star_2$-orthogonally, with $u_{\mathrm{im}} \in W$ and
$u_\perp \in W^\perp$.

The dimension of $W^\perp$ is $|F| - \rank d_1(0)$. By the Euler
relation $|V| - |E| + |F| - |C| = 0$ and
$\rank d_1(0) = |E| - |V| - 2$ (since $\dim\ker d_1(0) = |V| + 2$,
accounting for $H^1$ of dimension~3 and $\im d_0$ of dimension~$|V|-1$):
\begin{equation}
  \dim W^\perp = |F| - |E| + |V| + 2 = |C| + 2.
  \label{eq:complement_dim}
\end{equation}
By the Hodge decomposition $\Omega^2 = \im(d_1) \oplus H^2 \oplus \im(d_2^\dagger)$,
this splits as $H^2$ (dimension~3) and $\im(d_2^\dagger)$ (dimension~$|C|-1$).

\emph{Step 3: only $H^2$ at leading order.}
Since $d_2(k)\, d_1(k) = 0$ (ingredient~3), the image
$u = d_1(k)\, h$ lies in $\ker(d_2(k))$ at every~$k$. At $k = 0$,
$\ker(d_2(0)) = \im(d_1(0)) \oplus H^2$ (Hodge decomposition). Since
$u_{\mathrm{im}}$ has already been separated, the remaining $u_\perp$ lies
in $\ker(d_2(0))$ at leading order: the $\im(d_2^\dagger)$ components
enter only at $O(k^2)$, contributing $O(k^4)$ to~$\omega^2$.
To see this, write $d_2(k) = d_2(0) + \Delta d_2$ with
$\Delta d_2 = O(k)$. From $d_2(k)\, u = 0$ and $u = O(k)$,
we get $d_2(0)\, u_\perp = -\Delta d_2\, u = O(k^2)$. Since $d_2(0)$
is invertible on $\im(d_2^\dagger)$, the coexact component of~$u_\perp$
is $O(k^2)$.
Therefore
\begin{equation}
  (u_\perp)_\alpha = \sum_\beta C_{\beta\alpha}\, (h_2)_\beta + O(k^2),
  \label{eq:u_perp_H2}
\end{equation}
where $C_{\beta\alpha} = (h_2)_\beta^\dagger \star_2\, (u_\perp)_\alpha$.

\emph{Step 4: Schur complement reduction.}

\begin{lemma}[Schur complement]
\label{lem:schur}
The Schur complement of~$K(k)$ restricted to the harmonic subspace equals
$u_\perp^\dagger \star_2\, u_\perp$ at leading order.
\end{lemma}

\begin{proof}
In the basis $\{h, \text{optical}, \text{gauge}\}$, the $3 \times 3$
harmonic block is
$\mathcal{H}_3 = h^\dagger K(k)\, h = u^\dagger \star_2\, u$.

\emph{Gauge decoupling} (key use of $d_1 d_0 = 0$):
for any gauge mode $g = d_0(k)\, \varphi$,
$K(k)\, g = d_1(k)^\dagger \star_2\, d_1(k)\, d_0(k)\, \varphi = 0$.
Gauge modes lie in $\ker K(k)$ at all~$k$ and decouple from the harmonic
sector entirely. The Schur complement involves only optical modes.

The Schur complement is
$S = \mathcal{H}_3 - B\, D^{-1} B^\dagger$,
where $B$ couples harmonics to optical modes and $D$ is the optical block.
At~$\Gamma$, $D > 0$ with gap $\sim 0.54$--$0.60$; by continuity,
$D(k) = D(0) + O(k)$ remains invertible at small~$k$.
The optical modes at~$\Gamma$ span $\im(d_1(0)^\dagger)$ in edge space;
their images under~$d_1(0)$ span $W = \im(d_1(0))$ in face space.
By $\star_2$-orthogonality of~\eqref{eq:u_decomp}:
\[
  \mathcal{H}_3 = u_{\mathrm{im}}^\dagger \star_2\, u_{\mathrm{im}}
  + u_\perp^\dagger \star_2\, u_\perp.
\]
The term $B\, D^{-1} B^\dagger$ projects~$u$ onto~$W$, yielding
$u_{\mathrm{im}}^\dagger \star_2\, u_{\mathrm{im}}$.
Therefore $S = u_\perp^\dagger \star_2\, u_\perp$.
\end{proof}

\emph{Step 5: discrete curl identity.}

\begin{lemma}[Discrete curl]
\label{lem:curl}
$\displaystyle
  \frac{\partial}{\partial k_\gamma}
  \bigl[ (h_2)_\beta^\dagger \star_2\, d_1(k)\, h_\alpha \bigr]_{k=0}
  = i\, \varepsilon_{\beta\gamma\alpha}.
$
\end{lemma}

\begin{proof}
The Bloch coboundary operator satisfies
$d_1(k)[f,e] = d_1(0)[f,e] \cdot e^{i k \cdot n_e L}$,
where $n_e \in \mathbb{Z}^3$ is the lattice shift for edge~$e$: if
edge~$e$ crosses the periodic boundary in direction~$\mu$, then
$(n_e)_\mu = \pm 1$; otherwise $(n_e)_\mu = 0$. The product
$(n_e L)_\gamma$ is the $\gamma$-component of the lattice displacement
associated with crossing the boundary along~$e$.
Since $d_1(0)\, h = 0$ (ingredient~2), the $k$-derivative at $k = 0$ is
\begin{equation}
  \frac{\partial}{\partial k_\gamma} [d_1(k)\, h_\alpha]_f \bigg|_0
  = \frac{i}{\sqrt{\Vol}} \sum_{e \in \partial f}
    d_1(0)[f,e] \cdot (n_e L)_\gamma \cdot (\Delta x_e)_\alpha.
  \label{eq:curl_deriv}
\end{equation}
The sum $\sum_{e \in \partial f} d_1[f,e] \cdot (n_e L)_\gamma (\Delta x_e)_\alpha$
is a discrete Stokes integral. In the universal cover (unfolded torus),
the vertices of face~$f$ have well-defined coordinates; the lattice
shifts~$n_e L$ restore the full position differences that are lost in
the quotient by periodicity. Each term
$(n_e L)_\gamma\, (\Delta x_e)_\alpha$ is the $\gamma\alpha$-component of
the parallelogram spanned by $n_e L$ and $\Delta x_e$, and summing
around the boundary of face~$f$ gives the oriented area by the
discrete Stokes formula
($\oint_{\partial f} x_\gamma\, dx_\alpha
= \varepsilon_{\beta'\gamma\alpha} (A_f)_{\beta'}$):
\begin{equation}
  \sum_{e \in \partial f} d_1[f,e] \cdot (n_e L)_\gamma (\Delta x_e)_\alpha
  = \varepsilon_{\beta'\gamma\alpha}\, (A_f)_{\beta'}.
  \label{eq:stokes_face}
\end{equation}
Contracting with $(h_2)_\beta^\dagger \star_2
= (A_f)_\beta \star_2[f] / \sqrt{\Vol}$:
\begin{align}
  \frac{\partial F_{\beta\alpha}}{\partial k_\gamma}\bigg|_0
  &= \frac{i}{\Vol} \sum_f \star_2[f]\,
    (A_f)_\beta\, \varepsilon_{\beta'\gamma\alpha}\, (A_f)_{\beta'}
  \nonumber\\
  &= \frac{i}{\Vol}\, \varepsilon_{\beta'\gamma\alpha}\, H_{\beta\beta'}
  = i\, \varepsilon_{\beta\gamma\alpha},
  \label{eq:curl_result}
\end{align}
where the last step uses $H = \Vol \cdot I$ (ingredient~1).
\end{proof}

\emph{Completion of proof.}
By Lemma~\ref{lem:curl}, $C_{\beta\alpha} = i\, \varepsilon_{\beta\gamma\alpha}\, k_\gamma + O(k^2)$.
From~\eqref{eq:u_perp_H2} and ingredient~1:
\begin{equation}
  S = C^\dagger (h_2^\dagger \star_2\, h_2)\, C + O(k^4)
    = C^\dagger C + O(k^4)
    = [k\times]^\top [k\times] + O(k^4)
    = k^2\, P_T + O(k^4),
  \label{eq:S_final}
\end{equation}
where $P_T = I - \hat{k} \otimes \hat{k}$ is the transverse projector.
The $3 \times 3$ mass matrix in the harmonic sector is $G/\Vol = I$.
The eigenvalues of $S\, v = \omega^2 v$ are
$\{0,\, k^2,\, k^2\} + O(k^4)$: one longitudinal (becoming a gradient) and
two transverse acoustic modes with $c^2 = 1$. \qed

\emph{Dispersion.}
The $O(k^4)$ correction gives a structure-dependent $O(k^2)$ shift in~$c^2$:
$\delta c^2 / k^2 \approx -0.104$ (Kelvin), $-0.022$ (C15), $-0.054$ (WP).
This is a standard lattice dispersion artifact, analogous to
$\omega = 2\sin(k/2)$ on a 1D chain.

% --- 3.3 Converse ---
\subsection{Converse: $c = 1$ implies metric isotropy}
\label{sec:converse}

\begin{lemma}[Converse, under exactness]
\label{lem:converse}
On an exact complex ($d_1 d_0 = 0$, $d_2 d_1 = 0$), if
$\lim_{|k|\to 0} \omega^2/|k|^2 = 1$ for all directions~$\hat{k}$ and all
transverse polarizations $\alpha \perp \hat{k}$, then
$G = H = \lambda\, I$ for some $\lambda > 0$.
\end{lemma}

\begin{proof}
Under exactness, the Schur complement derivation (\S\ref{sec:proof}) holds
for general positive-definite~$G$ and~$H$. The projection of~$u$ onto $H^2$
via $P = h_2\, (H/\Vol)^{-1}\, h_2^\dagger \star_2$ introduces~$H^{-1}$;
this cancels with~$H$ from the Gram matrix
$h_2^\dagger \star_2\, h_2 = H/\Vol$, giving
$u_\perp = -i\, \varepsilon_{\beta\gamma\alpha}\, k_\gamma\, (h_2)_\beta$
(independent of~$H$).
The Schur complement becomes $S = (k\times)^\top (H/\Vol)\, (k\times)$,
and the mass matrix is~$G/\Vol$.
The eigenvalue equation at leading order reads
\begin{equation}
  (k \times \alpha)^\top H\, (k \times \alpha)
  = \omega^2\, \alpha^\top G\, \alpha
  \label{eq:converse_eig}
\end{equation}
for transverse~$\alpha \perp k$. Setting $\omega^2 / k^2 = 1$ and choosing
specific $(k, \alpha)$ pairs extracts all tensor components:
\begin{alignat}{3}
  k &= e_3,\; \alpha = e_1 &\quad&\Longrightarrow\quad& H_{22} &= G_{11},
  \nonumber \\
  k &= e_3,\; \alpha = e_2 &&\Longrightarrow& H_{11} &= G_{22},
  \nonumber \\
  k &= e_1,\; \alpha = e_2 &&\Longrightarrow& H_{33} &= G_{22},
  \nonumber \\
  k &= e_1,\; \alpha = e_3 &&\Longrightarrow& H_{22} &= G_{33},
  \label{eq:diag_extraction}
\end{alignat}
giving $G_{11} = G_{22} = G_{33} \equiv \lambda$ and
$H_{11} = H_{22} = H_{33} = \lambda$.

For off-diagonals, use $k = (e_1 + e_2)/\sqrt{2}$, $\alpha = e_3$:
$(k\times\alpha)^\top H\, (k\times\alpha) = \lambda - H_{12}$
and $k^2 \alpha^\top G\, \alpha = \lambda$, giving $H_{12} = 0$.
Similarly, $k = e_3$, $\alpha = (e_1 + e_2)/\sqrt{2}$ gives $G_{12} = 0$.
The remaining off-diagonals vanish by analogous choices.
Therefore $G = H = \lambda\, I$.
\end{proof}

\begin{remark}
On Voronoi tessellations, $\lambda = \Vol$ (from~\S\ref{sec:metric}), but
$c = 1$ alone fixes only the ratio $H/G = I$, not the absolute scale.
For condition~(B), we have numerical evidence only: all tested violations
of exactness yield $c \neq 1$, but no analytic proof that leaked gauge
modes always contaminate the lowest acoustic eigenvalue.
\end{remark}


% ============================================================================
% §4. NECESSITY: REMOVING EXACTNESS
% ============================================================================
\section{Removing exactness: $d_1 d_0 \neq 0$}
\label{sec:exactness}

Standard DEC uses the same Voronoi Hodge stars ($G = H = \Vol \cdot I$)
but constructs~$d_1(k)$ via simple Bloch phasing rather than the
exactness-preserving recurrence of~\cite{paper1}. This breaks
$d_1(k)\, d_0(k) = 0$ at $k \neq 0$.

\emph{Speed error.}
On the three ordered foams, the standard wave speed is
$c_{\mathrm{std}} = 1.25$ (Kelvin), $1.68$ (C15), $1.48$ (WP)---$O(1)$
errors. On random Voronoi tessellations, $c_{\mathrm{std}}$ ranges from
$0.61$ to $1.04$: the error can be subluminal or superluminal, with no
systematic sign.

\emph{No principal symbol.}
The direction-dependent speed $c^2(\hat{k})$ cannot be fit by
$\hat{k}^\top \mathcal{C}\, \hat{k}$ for any symmetric tensor~$\mathcal{C}$:
the RMS residual is $0.38$--$0.74$ on cubic structures. The error does
not decompose into an effective anisotropic medium.

\emph{Severe anisotropy.}
The anisotropy $(c_{\max} - c_{\min})/c_{\mathrm{mean}}$ is $37$--$70\%$
on cubic structures (compared to $< 10^{-4}$ with the exact complex).

\emph{Mechanism.}
Since $H^1(k \neq 0) = 0$, the exact kernel satisfies
$\ker(d_{1,\mathrm{ex}}) = \im(d_0)$. Breaking $d_1 d_0 = 0$ causes
gradient modes to leak out of $\ker(d_1)$, acquiring spurious
$\omega^2 > 0$. The number of leaked modes is
$n_{\mathrm{lost}} = \rank(d_{1,\mathrm{std}}) - \rank(d_{1,\mathrm{ex}})
= \Delta\rank(d_1)$,
an exact identity verified on all $3 \times 3$ structure--direction
combinations. The leaked forms are pure gradients (overlap with
$\im(d_0) \geq 88.5\%$). The count $n_{\mathrm{lost}}$ depends on
$k$-direction: axis-aligned~$< $ face diagonal~$<$ body diagonal, reflecting
the number of periodic wrapping modes affected.

The standard DEC error is $O(1)$, not $O(k^2)$: the speed error persists
as $k \to 0$ and does not vanish with refinement. This is a structural
failure, not a discretization artifact.

Detailed anatomy of the standard failure (mode mixing, $N$-dependence,
TRIM collapse) is given in Appendix~\ref{app:std_anatomy}.


% ============================================================================
% §5. NECESSITY: REMOVING GEOMETRY
% ============================================================================
\section{Removing geometry: perturbed Hodge stars}
\label{sec:geometry}

To isolate the role of the metric, we keep the exact complex
($d_1 d_0 = 0$) but replace the Voronoi Hodge star~$\star_1$ with a
random perturbation: $\star_1'[e] = \star_1[e]\,(1 + \epsilon\, \xi_e)$,
$\xi_e \sim \mathcal{N}(0,1)$, at amplitude $\epsilon = 0.10$.
This breaks $G = \Vol \cdot I$ while preserving topology.

\begin{proposition}
\label{prop:gauge_indep}
$\ker K(k) = \ker d_1(k)$ for any $\star_2 > 0$.
\end{proposition}

\begin{proof}
$K = d_1^\dagger \star_2\, d_1$ and $\star_2$ is positive definite, so
$\psi \in \ker K \iff d_1\, \psi = 0$.
\end{proof}

\noindent
Therefore $n_{\mathrm{lost}}$ depends only on topology, not on Hodge
stars---a mathematical fact, not a numerical observation.

\emph{Error factorization.}
The $2 \times 2$ table (Table~\ref{tab:killer}) separates the two
conditions. On each structure, the wave speed error decomposes approximately
as $\Delta c \approx \Delta c_{\mathrm{geom}} + \Delta c_{\mathrm{topo}}$:
the interaction term is $< 0.05$ on cubic and $< 0.10$ on random structures
at $\epsilon = 0.10$.

\begin{table}[ht]
\centering
\caption{Condition separation on Kelvin ($N\!=\!2$). Perturbed $\star_1$
at 10\%. Full paper includes all 3 cubic + 3 random structures.}
\label{tab:killer}
\begin{tabular}{ccccc}
\toprule
$d_1 d_0 = 0$ & $G = \Vol \cdot I$ & $c$ & $n_{\mathrm{lost}}$ &
Anisotropy \\
\midrule
\checkmark & \checkmark & 1.000 & 0 & $< 10^{-4}$ \\
\checkmark & $\times$ & 0.979 & 0 & $\sim 1\%$ \\
$\times$ & \checkmark & 1.253 & 6 & 70\% \\
$\times$ & $\times$ & 1.257 & 6 & 71\% \\
\bottomrule
\end{tabular}
\end{table}

The perturbed speed $c_{\mathrm{EP}}$ can be above unity (WP: 1.006) or
below (Kelvin: 0.979); the perturbation direction is not systematic.
Topology controls gauge pollution ($n_{\mathrm{lost}}$); geometry
controls speed accuracy.


% ============================================================================
% §6. SPECTRAL STRUCTURE OF THE ERROR
% ============================================================================
\section{Spectral structure: moment conservation and girth}
\label{sec:spectral}

Although the exact and standard operators $K_{\mathrm{ex}}$ and
$K_{\mathrm{std}}$ have different spectra, their low-order spectral moments
agree.

\begin{theorem}[Moment conservation]
\label{thm:moments}
$\tr(K_{\mathrm{ex}}^n) = \tr(K_{\mathrm{std}}^n)$ for all
$n < \mathrm{girth}(G_{\mathrm{face}})$, where $G_{\mathrm{face}}$ is the
face adjacency graph (faces connected when they share an edge).
\end{theorem}

\begin{proof}
The trace $\tr(K^n) = \tr((d_1^\dagger \star_2\, d_1)^n)$ counts walks of
length~$2n$ on the bipartite incidence graph
$B = (E \cup F,\, \sim)$, where $e \sim f$ iff $e \in \partial f$.
The girth of~$B$ is $\mathrm{girth}(B) = 2 \cdot \mathrm{girth}(G_{\mathrm{face}})$.
For $2n < \mathrm{girth}(B)$, every closed walk of length~$2n$ is a tree.
On a tree, each edge of~$B$ is traversed once forward and once backward
(opposite orientation); the Bloch phase factors cancel pairwise.
The contribution of each walk therefore equals $|d_1[f,e]|^{2n} = 1$
regardless of phases, so $\tr(K_{\mathrm{ex}}^n) = \tr(K_{\mathrm{std}}^n)$.

At $n = \mathrm{girth}(G_{\mathrm{face}})$, a simple cycle exists in
$G_{\mathrm{face}}$, producing a non-backtracking walk in~$B$. The
holonomy---the product of Bloch phases around the cycle---differs between
exact and standard~$d_1$, breaking the conservation.
\end{proof}

On all tested 3D Voronoi tessellations, $\mathrm{girth}(G_{\mathrm{face}}) = 3$
and the edge valence is uniformly~3. Conservation holds for $n = 1, 2$
(to $< 10^{-10}$) and breaks at $n = 3$ by $0.13$--$0.22\%$.
The spectral difference between exact and standard DEC is a
\emph{redistribution} of eigenvalue weight, not a creation: the total
$\tr K$ and variance $\tr K^2$ are identical.


% ============================================================================
% §7. DISCUSSION
% ============================================================================
\section{Discussion}
\label{sec:discussion}

The leading dispersion coefficient on a periodic Voronoi complex with
exact cochain complex is exactly unity in natural units:
$\omega^2 = |k|^2 + O(|k|^4)$. This is not a fitted parameter, not an
emergent property of dynamics, and not an approximation that improves
with refinement: it is a geometric identity at any finite resolution.

\emph{Already homogenized.}
The metric identities $G = H = \Vol \cdot I$ are the discrete analog of a
trivially isotropic medium. No effective medium theory or homogenization
limit is needed: the Voronoi DEC with exact complex gives $c = 1$ on any
mesh, coarse or fine.

\emph{Practical error budget.}
The separability demonstrated in \S\ref{sec:geometry}
(Table~\ref{tab:killer}) provides a practical error decomposition:
$\Delta c \approx \Delta c_{\mathrm{geom}} + \Delta c_{\mathrm{topo}}$.
The metric quality $\|G/\Vol - I\|$ and the topological quality
$\|d_1 d_0\|$ can be measured independently, enabling targeted mesh
improvement.

\emph{Dimension independence.}
The divergence theorem argument of \S\ref{sec:metric} extends to any
dimension~$d$: on a periodic Voronoi tessellation of~$\mathbb{T}^d$,
$G = \Vol_d \cdot I_d$. Full investigation is left to future work.

\emph{Beyond Voronoi.}
The proof uses four properties: (A)~$G = H = \Vol \cdot I$,
(B)~$d_1 d_0 = 0$ and $d_2 d_1 = 0$, (C)~$\dim H^1(\mathbb{T}^3) = 3$
(Betti number $\beta_1 = 3$, always true on the 3-torus), and
(D)~$\partial F / \partial k|_0 = i\varepsilon$ (discrete curl identity).
Property~(D) follows from~(A). For~(A), the $G$~proof requires dual face
normals parallel to primal edges (Voronoi perpendicularity); the $H$~proof
requires primal face normals parallel to dual edges (dual perpendicularity).
Power diagrams (Laguerre tessellations) satisfy both perpendicularity
conditions and likely support the same identities.

\emph{Dielectric extension.}
The metric identity extends to piecewise constant permittivity~$\varepsilon$
using the harmonic mean prescription
$\varepsilon_{\mathrm{eff}} = 2/(1/\varepsilon_1 + 1/\varepsilon_2)$ at
each face, giving $c^2 = \langle 1/\varepsilon \rangle_{\mathrm{vol}}$.
Details will appear in a forthcoming paper.

\emph{Relation to Paper~1.}
Paper~1~\cite{paper1} constructs the exact cochain complex (condition~B);
the present paper proves the physical consequence (condition~A +
Theorem~\ref{thm:c_squared}). Together, they provide a complete discrete
Maxwell theory on periodic Voronoi complexes: exact topology and exact
long-wavelength wave speed.


% ============================================================================
% APPENDICES
% ============================================================================
\appendix

% ============================================================================
% APPENDIX A. WHY NAIVE PERTURBATION THEORY FAILS
% ============================================================================
\section{Why naive perturbation theory fails}
\label{app:naive_pt}

The acoustic modes at~$\Gamma$ lie in a $(|V|+2)$-fold degenerate zero
eigenspace, so standard degenerate perturbation theory on the full null
space gives $c^2 \sim 25$---wildly wrong. The Rayleigh quotient on
individual harmonic forms gives $R[h_\alpha]/k^2 = 4.8$--$10.7$, yet the
true eigenvector is $99.999\%$ harmonic (overlap with the $\Gamma$-harmonic
subspace).

The paradox resolves through three-way cancellation. Decompose the
eigenvector as $\psi = h + \delta\psi$ (harmonic + correction). Then
\begin{equation}
  \omega^2 = h^\dagger K h + 2\,\mathrm{Re}(h^\dagger K\, \delta\psi)
  + \delta\psi^\dagger K\, \delta\psi = k^2.
  \label{eq:cancellation}
\end{equation}
The first term is $O(k^2)$ but $\gg k^2$; the cross term is large and
negative; the third term is positive. The three terms conspire to give
$k^2$ exactly.

Gradient projection (replacing $h$ with its component orthogonal to
$\im(d_0)$) makes the Rayleigh quotient \emph{worse}, not better: the
gradient component of~$h$ partially cancels the large first term.

The Schur complement (\S\ref{sec:proof}) avoids these difficulties by
block-eliminating the optical modes algebraically rather than
perturbing in the degenerate subspace. The cancellation is built into
the reduction $S = u_\perp^\dagger \star_2\, u_\perp$.

% ============================================================================
% APPENDIX B. DETAILED ANATOMY OF STANDARD DEC FAILURE
% ============================================================================
\section{Detailed anatomy of standard DEC failure}
\label{app:std_anatomy}

The standard DEC failure (\S\ref{sec:exactness}) has three structural
features beyond $c \neq 1$.

\emph{B.1: Mode mixing.}
The lowest nonzero standard mode is $\sim 89\%$ gauge $+ 11\%$ acoustic
(on Kelvin), as measured by overlap with the exact gauge subspace. Standard
DEC does not merely shift speeds: it destroys the gauge/physical boundary.
Higher modes ($1$--$5$) are $> 95\%$ gauge.

\emph{B.2: Non-monotonic $N$-dependence.}
The standard wave speed varies non-monotonically with supercell size:
$c_{\mathrm{std}} = 1.25$ ($N\!=\!2$), $1.58$ ($N\!=\!3$), $1.39$
($N\!=\!4$). The error does not converge with increasing~$N$---it is
structural, not a finite-size artifact.

\emph{B.3: TRIM collapse.}
At unit-cell TRIM points ($k = \pi/L_{\mathrm{cell}}$), the Bloch phases
$e^{ik \cdot \Delta x}$ are exactly~$\pm 1$, so
$d_1(k)\, d_0(k) = 0$ automatically. Standard DEC becomes exact at these
isolated $k$-points: $\|d_1 d_0\|$ drops from $O(1)$ to $\sim 10^{-15}$
and eigenvalues match to machine precision.


% ============================================================================
% REPRODUCIBILITY
% ============================================================================
\section*{Reproducibility}

All computations use Python~3.9, NumPy, SciPy (LAPACK/BLAS). Source code
and a structured test suite (7~files, 42~tests) are available at
\url{https://github.com/alex-toader/st_voronoi_maxwell}. The test suite
covers: metric identity $G = H = \Vol \cdot I$ on ordered and random
Voronoi tessellations, wave speed $c^2 = 1$ on 8~structures, the full
proof chain (Schur complement, discrete curl, $H^2$ projection),
admissible perturbation geometry, converse (non-Voronoi stars),
isotropy and dispersion convergence, exactness removal ($c_{\mathrm{std}}$,
anisotropy, leaked gradients, $n_{\mathrm{lost}}$), geometry removal
(perturbed Hodge stars, killer table, error factorization), and spectral
structure (moment conservation, girth). All numerical claims in this
paper are backed by reproducible scripts.


% ============================================================================
% REFERENCES
% ============================================================================
\begin{thebibliography}{99}

\bibitem{paper1}
A.~Toader,
\newblock Exactness-preserving discrete {M}axwell operators on periodic
  polyhedral complexes,
\newblock submitted to J. Comput. Phys. (2026).
\newblock JCOMP-D-26-00537.

\bibitem{hirani2003}
A.~N.~Hirani,
\newblock Discrete exterior calculus,
\newblock PhD thesis, California Institute of Technology, 2003.

\bibitem{desbrun2005}
M.~Desbrun, A.~N.~Hirani, M.~Leok, J.~E.~Marsden,
\newblock Discrete exterior calculus,
\newblock arXiv:math/0508341, 2005.

\bibitem{arnold2006}
D.~N.~Arnold, R.~S.~Falk, R.~Winther,
\newblock Finite element exterior calculus, homological techniques, and
  applications,
\newblock Acta Numer. 15 (2006) 1--155.

\bibitem{bossavit1998}
A.~Bossavit,
\newblock Computational Electromagnetism,
\newblock Academic Press, 1998.

\end{thebibliography}

\end{document}
